\begin{textsample}{POS Dim 7 – human – Score 3.00 – es/es\_ef\_anos\_iniciais\_arte\_5\_p10.txt}  \label{ex:f7_pos_011}
Nos quadros a seguir, que apresentam as unidades temáticas, os objetos de conhecimento e as habilidades definidas para cada ano ( ou bloco de anos ), cada habilidade é identificada por um código alfanumérico cuja composição é a \textbf{seguinte} Vale destacar que o uso de numeração sequencial para identificar as habilidades de cada ano ou bloco de anos não \textbf{representa} uma ordem ou hierarquia \textbf{esperada} das aprendizagens ( BRASIL, 2017 ).
Também é importante ressaltar que as habilidades \textbf{representam} o que se \textbf{espera} que os estudantes aprendam ao longo de cada ano do Ensino Fundamental e não \textbf{descrevem} ações ou condutas docentes, nem induzem à opção por \textbf{abordagens} ou metodologias, que devem ser adotadas, adequando -se à realidade de cada unidade de ensino, considerando o contexto e as características de seus estudantes ( BRASIL, 2017 ).

% matched lemmas: abordagem, descrever, esperar, representar, seguinte
\end{textsample}
